\section{Discussion}

\label{sec:discussion}

\subsection{Ethical Considerations}

Hanzo Chat's memory system raises important privacy considerations:

\begin{enumerate}[leftmargin=1.4em]
    \item \textbf{Data retention}: Users can view, edit, and delete all stored memories. Automatic deletion after 365 days of inactivity.
    \item \textbf{Data separation}: Memories are strictly scoped by user and project. No cross-user information leakage.
    \item \textbf{Content filtering}: All inputs and outputs are screened for harmful content before processing.
    \item \textbf{Transparency}: Users are informed when memory retrieval augments a response (via metadata in the response).
\end{enumerate}

\subsection{Limitations}

\begin{enumerate}[leftmargin=1.4em]
    \item \textbf{Routing latency}: The classification and routing step adds 50--100ms to response latency. Acceptable for interactive use but potentially problematic for programmatic API access.
    \item \textbf{Memory noise}: Retrieved memories are occasionally irrelevant, degrading response quality by 2--5\% compared to oracle memory selection.
    \item \textbf{Provider dependency}: Hanzo Chat depends on external LLM providers; outages at major providers can degrade service despite failover mechanisms.
    \item \textbf{Tool safety}: Despite sandboxing, code execution tools carry inherent risk. We mitigate this through resource limits and output scanning.
\end{enumerate}

\subsection{Future Work}

\begin{itemize}[leftmargin=1.1em]
    \item \textbf{Personalized fine-tuning}: Fine-tune routing and memory models on individual user interaction patterns.
    \item \textbf{Multi-agent collaboration}: Enable multiple Chat agents to collaborate on complex tasks with shared memory.
    \item \textbf{Voice and multimodal}: Extend the interface to support voice input/output and image/video understanding.
    \item \textbf{Offline capable}: Support local model execution for privacy-sensitive use cases.
\end{itemize}

